\scanpage{47}
\mspage{46}
\deleted{
- \mspage{IV-1007} -

plication b)$\Rightarrow$c)\inserted{sans hypothèse de séparation}. Nous ignorons s'il en est encore ainsi (sans hypothèse de séparation) pour $n$ quelconque. Rappelons que, sans hypothèse de séparation, l'implication c)$\Rightarrow$b) n'est plus valable\inserted{même pour $n=1$} (0,15.1.12 (iii)).

\uline{19.6. Suites régulières relativement à un module \deleted{quotient}\inserted{filtré}}.

(19.6.1) Soient $A$ un anneau, $\underline{f}=(f_1,\ldots,f_n)$ une suite\inserted{finie} d'éléments de $A$, $\underline{J}=f_1A+\ldots+f_nA$ l'idéal qu'elle engendre, et considérons une suite exacte de $A$-modules
\diagram{\begin{tikzcd}
0\arrow[r]&R\arrow[r,"j"]&N\arrow[r,"p"]&M\arrow[r]&0
\end{tikzcd}}
où l'on suppose que $N$ est muni d'une filtration décroissante $(N_k)$ formée de sous-$A$-modules, avec $N_0=N$. On pose $R_k=R\cap N_k$ (filtration induite par $(N_k)$), $M_k=p(N_k)$ (filtration quotient de $(N_k)$), et on désigne par $\operatorname{gr}_\bullet(N)$, $\operatorname{gr}_\bullet(R)$ et $\operatorname{gr}_\bullet(M)$ les $A$-modules gradués associés à ces trois filtrations. On pose d'autre part pour tout $k$,
\[
N'_k=N_k+\underline{J}N_{k-1}+\ldots+\underline{J}^{k-1}N_1+\underline{J}^kN_0
\]
de sorte que \deleted{$\unsure{M'_k=p(N'_k)}$} $M'_k=p(N'_k)=M_k+\underline{J}M_{k-1}+\ldots+\underline{J}^{k-1}M_1+\underline{J}^kM_0$ ; on pose d'autre part $R'_k=R\cap N'_k$ (filtration \uline{induite} par $(N'_k)$), et on note \deleted{\unsure{par}} $\operatorname{gr}'_\bullet(N)$, $\operatorname{gr}'_\bullet(M)$ et $\operatorname{gr}'_\bullet(R)$ les $A$-modules gradués associés à ces trois filtrations ; on a donc un diagramme commutatif de suites exactes

(19.6.1.1)
\diagram{\begin{tikzcd}[column sep=small]
0\arrow[r]&\operatorname{gr}_\bullet(R)\arrow[r,"\operatorname{gr}(j)"]\arrow[d]&\operatorname{gr}_\bullet(N)\arrow[r,"\operatorname{gr}(p)"]\arrow[d]&\operatorname{gr}_\bullet(M)\arrow[r]\arrow[d]&0\\
0\arrow[r]&\operatorname{gr}'_\bullet(R)\arrow[r,"\operatorname{gr}'(j)"']&\operatorname{gr}'_\bullet(N)\arrow[r,"\operatorname{gr}'(p)"']&\operatorname{gr}'_\bullet(M)\arrow[r]&0
\end{tikzcd}}
\inserted{où les flèches verticales sont les homomorphismes canoniques d'un module gradué associé à une filtration dans le module gradué associé à une filtration moins forte.}

On notera que si l'on pose
\[
R''_k=R_k+\underline{J}R_{k-1}+\ldots+\underline{J}^{k-1}R_1+\underline{J}^kR_0
\]
on a évidemment $R''_k\subset R'_k$, mais les filtrations $(R''_k)$ et $(R'_k)$ sont en général distinctes ; on notera $\operatorname{gr}''_\bullet(R)$ le $A$-module gradué associé à la filtration $(R''_k)$.

(19.6.2) On a défini en (\uline{0},15.1.5.2) les homomorphismes surjectifs gradués de degré 0
\[
\psi_N:(\operatorname{gr}_\bullet(N)\otimes_A(A/\underline{J}))[T_1,\ldots,T_n]\longrightarrow\operatorname{gr}'_\bullet(N)
\]
}
\figure{\begin{tikzpicture}[x=.32cm,y=.32cm]
\draw (2.7,8.7) .. controls (2.1,9) and (4.4,15.9) .. (7.6,19.7);
\draw (3.8,2.7) .. controls (4.5,2.6) and (10.1,17.3) .. (12.5,19.7);
\draw (7.5,2.2) .. controls (8.4,.8) and (10.7,9.3) .. (16.5,12.7);
\draw (1.5,16.5) .. controls (1.9,17) and (2.7,17.1) .. (2.5,17.3);
\end{tikzpicture}}
